% Metódy inžinierskej práce, ak. rok 2026/27
% Spresnenie zamerania projektu

\documentclass[11pt,a4paper,slovak]{article}

\usepackage[slovak]{babel}
\usepackage[IL2]{fontenc}
\usepackage[utf8]{inputenc}
\usepackage[margin=2.5cm]{geometry}
\usepackage{url}

\title{Spresnenie zamerania projektu}

\author{Bohdan Kochkin \and Oleksandr Komarenko \and Dmytro Kolodeznyi \and Daniel Kollár\\[6pt]
	{\small Slovenská technická univerzita v Bratislave}\\
	{\small Fakulta informatiky a informačných technológií}\\[4pt]
	{\small Cvičiaci: Ing. Oliver Udvardi}
	}

\date{\small 3. október 2026}

\begin{document}

\maketitle

\section*{Názov projektu}

Detekcia podozrivých kryptomenových transakcií pomocou grafových neurónových
sietí

\section*{Akronym projektu}

STOPA (\emph{Sledovanie Transakčných Operácií a Podozrivých Aktivít})

\section*{Anotácia}

Transakcie v~blockchaine sú verejné, osoby za adresami však zostávajú anonymné.
Páchatelia preháňajú prostriedky z~trestnej činnosti cez dlhé reťazce prevodov,
aby zahladili ich pôvod. Samotná transakcia zriedka vyzerá podozrivo, podozrivou
ju robí okolie: odkiaľ prostriedky prišli a kam smerujú ďalej. Túto štruktúru
nezohľadňujú riešenia založené na znakoch samostatnej transakcie. A hoci sa
grafové metódy aktívne skúmajú, ich výhoda oproti jednoduchším modelom nie je
vždy potvrdená.

V~projekte STOPA skúmame, či grafové neurónové siete skutočne lepšie odhalia
podozrivé transakcie než klasické modely. Transakcie si predstavujeme ako graf,
kde uzol je transakcia a hrana tok peňazí medzi nimi. Klasické modely ---
logistická regresia, random forest, gradient boosting --- pracujú len
s~tabuľkovými dátami a mnohé súvislosti im uniknú, zatiaľ čo graf ich dokáže
zachytiť. Testujeme preto aj hybrid, ktorý spája grafové príznaky s~klasickým
klasifikátorom, a overujeme, či prináša reálne zlepšenie.

Vyhodnotenie prebieha na verejnom anotovanom datasete Elliptic s~viac ako
200~000 bitcoinovými transakciami. Dáta delíme podľa časových období, čím
minimalizujeme únik informácií z~budúcnosti a sledujeme, ako sa výkonnosť modelu
mení v~čase. Vzhľadom na nevyváženosť tried sa zameriame na úplnosť nelegálnej
triedy aspoň na úrovni základného modelu a na zvýšenie F1 skóre hybridu aspoň
o~5 percentuálnych bodov, prípadne so zdôvodnením, ak sa to nepodarí. Súčasťou
hodnotenia bude plocha pod precision-recall krivkou (PR-AUC) a sledovanie
poklesu kvality v~čase.

Výsledky sú využiteľné pre finančnú spravodajskú jednotku, orgány činné
v~trestnom konaní a poskytovateľov služieb kryptoaktív podliehajúcich nariadeniu
MiCA. Venujeme sa aj obmedzeniu nesprávnych označení, keďže bezdôvodné
podozrenie voči osobe prináša reálne riziká. Výstupmi budú prototyp analytického
reťazca, porovnávacia štúdia metód a prehľadový článok.

\end{document}